\section{实验设计与探索过程}

\subsection{六组主实验}

表\ref{tab:experiment-matrix}不是简单的模型排行榜，而是对应四个连续问题。E1--E3 先建立中心训练基线；E4 检查扩大训练分布的效果；E5 改变 token 聚合方式；E6 组合局部卷积特征与 Transformer。所有模型使用同一数据划分、优化器和评估协议。

\begin{table}[H]
  \centering
  \caption{主实验矩阵及各实验回答的问题}
  \label{tab:experiment-matrix}
  \resizebox{\textwidth}{!}{%
    \begin{tabular}{lllll}
      \toprule
      编号 & 模型 & 训练分布 & 关键设置 & 主要问题 \\
      \midrule
      E1 & MLP & Center & 展平输入 & 缺少空间结构时位置退化有多明显？ \\
      E2 & CNN & Center & 局部卷积与池化 & 权值共享能否自然抵抗大幅平移？ \\
      E3 & ViT-AbsPos & Center & Learnable Pos + CLS & 全局交互是否仍会学习中心偏置？ \\
      E4 & ViT-Aug & Shift+Rotation & CLS + Random Erasing & 只扩大训练分布能改善多少？ \\
      E5 & ViT-MeanPool & Shift+Rotation & Mean Pooling & 聚合方式是否改善稳定性？ \\
      E6 & HybridConv-ViT & Shift+Rotation & Conv Stem + Mean & 局部归纳偏置能否继续提升？ \\
      \bottomrule
    \end{tabular}%
  }
\end{table}

\subsection{阶段一：中心基线是否可靠}

最先需要验证的不是某个改进结构，而是原始中心测试是否掩盖问题。E1--E3 的训练图像始终位于中心，测试时逐步扩大平移幅度。如果 Center Accuracy 保持较高而 Grid Accuracy 大幅下降，就能直接说明“会分类”和“能适应位置变化”是两个不同目标。

三种基线覆盖了不同空间假设：MLP 基本不保留二维结构，CNN 具有局部共享，ViT-AbsPos 允许全局 token 交互但显式编码绝对坐标。它们之间的比较用于描述失效形态，而不是仅按中心准确率确定胜负。

\subsection{阶段二：先处理数据分布}

在确认中心偏置后，E4 保持 ViT-AbsPos 的结构和参数量不变，将训练输入改为随机平移与 $\pm45^\circ$ 旋转，并加入随机擦除。这一阶段检验最直接的工程假设：模型只有在训练期看过足够广的位置和角度，才可能在测试时形成平坦响应。

由于 E4 同时包含平移、旋转和擦除，它不是严格的单因素消融。选择这种组合是为了先建立一个有效的增强方案，再根据结果决定哪些因素值得拆分。固定角度扫描是后续补充的重要环节：它避免随机 Rotation Accuracy 把极端角度的下降平均掉。

\subsection{阶段三：聚合方式是否重要}

E5 在 E4 基础上只把 CLS token 改为 Mean Pooling。两者模型参数量和训练增强基本一致，适合观察聚合方式带来的增量。如果 Grid Accuracy、角度平均或最差位置进一步改善，可以提出“让所有 patch 直接参与表示有助于稳定性”的解释；如果差异很小，则说明主要收益可能已经来自训练分布覆盖。

\subsection{阶段四：加入局部视觉先验}

E6 在 E5 前加入两层卷积 stem。该阶段关注的是小型灰度图像上的局部模式：当主体被平移、旋转或截断时，袖口、鞋底、包带等局部边缘仍可能存在。卷积权值共享可以先提取这些重复出现的局部特征，再由 Transformer 整合全局关系。

E6 是组合方案而非严格的 stem 单变量实验。因此，若 Hybrid 表现最好，合理表述是“卷积 stem、Mean Pooling 与增强的组合在当前预算下最稳定”，而不是断言提升全部由卷积产生。

\subsection{评价证据链}

每个研究问题由不同证据回答：

\begin{itemize}
  \item Center、Small Shift、Large Shift 和 Robust Drop 描述从原分布到随机平移分布的总体变化；
  \item 49 点网格热力图揭示高准确率是否只集中在画布中央，并给出最差位置；
  \item Rotation、Shift+Rotation 和固定角度扇形图描述朝向变化与复合扰动；
  \item 逐类准确率、混淆矩阵和错误样本解释平均值背后的类别与可见性差异；
  \item attention rollout 用于观察响应区域是否随前景移动，但不作为模型决策的因果证明。
\end{itemize}

这种设计将“现象是否存在”“改进是否有效”和“仍然在哪里失败”分开处理，避免只用一项最高 Accuracy 概括整个项目。
